\chapter*{Conclusion générale et perspectives}
\addcontentsline{toc}{chapter}{Conclusion générale et perspectives}

\section*{Perspectives}
\addcontentsline{toc}{section}{Perspectives}

Pour la suite du projet, plusieurs perspectives d'évolution peuvent être envisagées afin d'enrichir SIRH-Doc et de l'adapter à des usages administratifs plus larges. Une première évolution concernerait l'intégration d'un circuit de validation documentaire. Dans ce scénario, un document pourrait être préparé par un utilisateur, vérifié par un responsable, puis validé avant son archivage définitif. Cette fonctionnalité renforcerait la fiabilité des pièces produites, surtout pour les documents sensibles comme les contrats, les attestations ou les décisions administratives.

Une autre perspective importante serait l'ajout de la signature électronique. Elle permettrait de donner une valeur plus forte aux documents générés et de réduire davantage les échanges papier. Cette évolution pourrait être complétée par l'envoi automatique des documents par courriel, notamment pour les vacataires, les enseignants ou les agents qui doivent recevoir rapidement une copie officielle.


\section*{Conclusion générale}
\addcontentsline{toc}{section}{Conclusion générale}

La réalisation de SIRH-Doc a constitué une expérience complète d'ingénierie logicielle. Le projet ne s'est pas limité à la création d'une interface de génération de documents ; il a demandé de comprendre un contexte métier réel, de tenir compte des contraintes d'un établissement universitaire et de concevoir une solution capable de rester adaptable d'une organisation à l'autre.

Tout au long du projet, plusieurs connaissances acquises durant la formation ont été mobilisées : la modélisation UML, la conception orientée objet, l'architecture logicielle, le développement web, la gestion des bases de données et l'organisation du travail par sprints. Ces éléments ont été appliqués dans un cadre concret, avec des choix techniques qui devaient répondre à des besoins précis : isoler les données, gérer les rôles, prendre en compte l'année universitaire, configurer les modules et assurer la traçabilité des documents produits. Ce travail nous a aussi permis de mieux comprendre la manière dont un projet est conduit dans un environnement professionnel, avec ses contraintes, ses échanges et ses ajustements progressifs.

Le projet a également permis d'approfondir de nouvelles compétences. La mise en place d'un moteur basé sur les métadonnées, la génération documentaire à partir de templates, la gestion des variables tableau et l'archivage des documents ont posé des problèmes qui dépassent les exercices classiques de développement. L'utilisation de Tiptap nous a notamment fait découvrir plus concrètement le monde de l'open source et la manière d'adapter une bibliothèque existante à un besoin métier précis.

Malgré les contraintes de temps liées au cadre d'un projet de fin d'études, le travail réalisé a permis d'aboutir à une plateforme cohérente et fonctionnelle. Les quatre sprints ont progressivement construit le socle de sécurité, le moteur de données, le studio documentaire puis le module de production et d'archivage. Cette progression a donné au projet une structure lisible et a facilité la validation des fonctionnalités au fur et à mesure.

En définitive, SIRH-Doc représente une base solide pour accompagner la dématérialisation des documents RH dans un établissement éducatif. Le système reste perfectible, comme tout produit appelé à évoluer, mais ses fondations techniques et fonctionnelles permettent déjà d'envisager des extensions sans remettre en cause l'ensemble de l'architecture. Ce projet a ainsi été l'occasion de relier les acquis académiques à une problématique réelle, tout en développant une solution utile, structurée et ouverte à de futures améliorations.
